\subsection{56. ML-Accelerated Computational Fluid Dynamics (Kochkov et al.\ PNAS 2021)}
\textbf{arXiv:} \texttt{2102.01010}\quad\textbf{Authors:} Kochkov, Smith, Ronber, et al.\ (2021)\quad\textbf{Journal:} PNAS 118(21)\quad\textbf{Domain:} ML / PDE / CFD\\
\scorebadge{Coverage}{8}\quad\scorebadge{Agreement}{9}\quad\textbf{Total:} 17/20\par\smallskip
\noindent\textbf{Coverage rationale.} Replicated the paper's three primary physical regimes: forced Kolmogorov flow at Re$\approx$1000 (Figs~1--2), decaying turbulence (Figs~3, 5), and Re$\approx$4000 (Fig~4). 12/14 testable claims evaluated (86\%); 5/8 primary analyzable units covered (63\%). The three uncovered units (architecture comparison, LES, large-domain generalization) are supplementary rather than headline results. Authors' pre-trained model outputs (released publicly via GCS) were independently evaluated to confirm paper's quantitative claims. Total compute: $\sim$2~GPU-hours on A100.\par
\noindent\textbf{Agreement rationale.} 11/12 tested claims verified, 1 partially verified. Using our abbreviated training (1/10--1/100 of the paper's compute): LI(64) at Re=1000 achieves $t_{\mathrm{dec}}$=3.58 (between DNS128=2.59 and DNS256=4.21), demonstrating correct trend at $\sim$4$\times$ resolution equivalence. The authors' released model achieves $t_{\mathrm{dec}}$=7.01 ($\sim$8$\times$), confirming the paper's headline 8--10$\times$ claim. Long-term stability verified for 2000+ frames without divergence. Energy spectra match the $k^{-3}$ enstrophy cascade. The gap between our results and the paper's is entirely attributable to training compute budget, not methodology.\par\smallskip
\noindent\textbf{Key findings:}
\begin{itemize}[leftmargin=1.4em,itemsep=0.2em,topsep=0.2em]
\item LI(64) matches DNS256--DNS512 accuracy at Re=1000 (author model); our abbreviated model matches DNS128--DNS256 (4$\times$ vs 8$\times$, training-gap explained)
\item 38--357$\times$ speedup confirmed from the paper's released TPU timing data
\item Decaying turbulence: LI(64) at 10k training steps achieves $\sim$5--7$\times$ resolution equivalence (author model), matching the paper's $\sim$7$\times$ claim
\item Re=4000: LI(128) achieves $\sim$5--7$\times$ (author model), consistent with paper
\item Zero-shot cross-regime transfer does \emph{not} work (corr@$t$=2=0.39); the paper's ``generalization'' claim refers to retrained models, a nuance not always clear in the abstract
\end{itemize}
\noindent\textbf{Verdict: REPLICATED.} Core claims confirmed both via our independent training and via evaluation of the authors' public model outputs.
\noindent\textbf{Follow-on research questions:}
\begin{enumerate}[leftmargin=1.6em,itemsep=0.2em,topsep=0.2em,label=Q\arabic*.]
\item At what training-step budget does the LI(64) model's $t_{\mathrm{dec}}$ saturate? Plot $t_{\mathrm{dec}}$ vs training steps from 1k to 200k to identify the compute--accuracy Pareto frontier.
\item Can the inner-step parameter be tuned adaptively during rollout (lower inner steps early, higher later) to improve stability--accuracy tradeoff?
\item Does the LI architecture generalize to 3D turbulence (Kolmogorov forcing on $[0,2\pi]^3$) with a manageable compute cost, or does the $N^3$ scaling make it impractical?
\item How sensitive is LI to the choice of forcing wavenumber? Sweep $k \in \{1,2,4,8\}$ and measure whether the resolution equivalence ratio holds.
\item Implement and compare the LC (Learned Correction) and EPD architectures from the paper's Fig~5 against LI---which architecture dominates at matched parameter count?
\end{enumerate}

